\section{Теоретические основы}

\subsection{Метаграфы как инструмент моделирования сложных систем}

Метаграфы были впервые предложены как обобщение направленных графов для моделирования систем с гиперсвязями. В отличие от классического графа, где дуга соединяет две вершины, в метаграфе дуга может соединять \textbf{множества вершин}, что позволяет естественным образом описывать отношения «многие ко многим». Это особенно важно в контексте информационной безопасности, где одно действие (например, запуск скрипта) может одновременно затрагивать несколько объектов: процесс, файл, сетевой сокет, учётную запись.

Формально метаграф задаётся как  
\[
G = (V, E),
\]  
где \(V\) — конечное множество вершин, \(E \subseteq 2^V \times 2^V\) — множество дуг. Каждая дуга \(e = (S, T) \in E\) интерпретируется как переход от множества источников \(S \subseteq V\) к множеству целей \(T \subseteq V\).

\subsection{Атрибутивные расширения}

Для применения в ИБ метаграф необходимо расширить атрибутами, отражающими семантические свойства компонентов системы. Вводятся функции атрибуции:
\[
A_V: V \to \mathcal{A}_V, \quad A_E: E \to \mathcal{A}_E,
\]  
где \(\mathcal{A}_V\) и \(\mathcal{A}_E\) — домены атрибутов.

Примеры атрибутов:
- Для вершины «пользователь»: роль (администратор, гость), уровень доверия, принадлежность к группе.
- Для вершины «процесс»: исполняемый файл, хеш, родительский процесс.
- Для дуги «доступ к файлу»: тип доступа (чтение/запись), протокол, временная метка, тактика MITRE ATT\&CK.

Такое расширение позволяет строить \textbf{контекстно-зависимые модели}, в которых не только структура, но и семантика играют ключевую роль в анализе.

\subsection{Связь с существующими фреймворками}

Предлагаемая модель органично интегрируется с:
- \textbf{MITRE ATT\&CK} — тактики и техники становятся атрибутами дуг;
- \textbf{STIX 2.1} — объекты STIX (indicator, attack-pattern, malware) могут быть отображены в вершины метаграфа;
- \textbf{Cyber Kill Chain} — этапы цепочки атаки соответствуют уровням модели.

Это обеспечивает совместимость с существующей экосистемой кибербезопасности и упрощает внедрение.